\subsection{Terminology and Etiological Frameworks}
Before surveying the theories proposed to explain cybersickness, it is useful to delimit the construct itself and to distinguish it from related forms of sickness and from perceptual phenomena that often co-occur with it.

Motion sickness is the broadest construct, encompassing adverse responses to both physical and visually specified motion \cite{lacknerMotionSicknessMore2014, keshavarzMotionSicknessCurrent2022}.
Within it, visually induced motion sickness (VIMS) designates sickness elicited primarily by visual motion, including cases where physical body motion is absent or inconsistent with the visual scene \cite{kennedyResearchVisuallyInduced2010, chaMotionSicknessDiagnostic}.
Cybersickness is most often treated as the VR-specific instance of VIMS; simulator sickness, gaming sickness, and related labels denote overlapping exposure contexts rather than distinct symptom categories \cite{cakirHandbookVirtualEnvironments2015, stanneyCybersicknessNotSimulator1997, rebenitschReviewCybersicknessApplications2016}.
This review adopts the above hierarchy when introducing constructs, but reports each included study using the terminology its authors chose, in order to avoid imposing a uniform vocabulary on a literature that does not yet share one.

Traditionally, these constructs have been measured using subjective questionnaires.
The dominant instrument is the Simulator Sickness Questionnaire (SSQ), which yields nausea, oculomotor, and disorientation subscales \cite{kennedySimulatorSicknessQuestionnaire1993}; it is frequently substituted by the single-item Fast Motion Sickness Scale (FMS) for fast, continuous, in-task ratings \cite{keshavarzValidatingEfficientMethod2011}.
Broader diagnostic criteria for motion sickness and VIMS extend these subscales to gastrointestinal, thermoregulatory, arousal, dizziness or vertigo, headache, and visual or oculomotor domains \cite{chaMotionSicknessDiagnostic}.

Several etiological frameworks have been proposed to explain cybersickness and the wider VIMS phenomenon; they are not mutually exclusive and often emphasise complementary mechanisms.

The most commonly used framework is the \textit{sensory conflict theory}, which attributes sickness to contradictory visual, vestibular, and proprioceptive signals, or between current sensory input and the prior exposure history encoded by the observer \cite{reasonMotionSickness1975, reasonMotionSicknessAdaptation1978, omanMotionSicknessSynthesis1990}.
Two refinements within this family further specify the nature of the conflict: subjective vertical mismatch frames it as a discrepancy between sensed and expected gravitational vertical \cite{blesMotionSicknessOnly1998}, while rest-frame conflict locates it in the observer's difficulty identifying a stable spatial reference \cite{protheroUnifiedApproachPresence2003}.
This family is particularly pertinent to VR, where HMDs routinely present visual self-motion while vestibular and proprioceptive signals indicate little or no corresponding physical movement.

A second framework, \textit{postural instability theory}, locates the cause of sickness not in sensory mismatch but in a sustained inability to maintain postural control under unfamiliar sensorimotor demands \cite{riccioEcologicalTheoryMotion1991, stoffregenPosturalInstabilityPrecedes1998}.
The theory is most directly testable in VR paradigms involving standing, walking, or large-field optic flow, where body sway and control dynamics can be measured objectively and have been reported to precede or accompany subjective sickness \cite{cortesEEGbasedExperimentVR2023a}.
A related strand of recent VR work returns to the role of vection itself, showing that sickness intensifies when visually induced self-motion violates the user's expected motion or control state \cite{teixeiraUnexpectedVectionExacerbates2022, teixeiraInvestigatingRelativeContributions2024}.

A further set of proposals address narrower mechanisms or different levels of explanation.
Vergence-accommodation conflict identifies a mismatch between binocular convergence and the fixed focal distance of stereoscopic displays as a source of oculomotor discomfort \cite{kramidaResolvingVergenceAccommodationConflict2016}, while oculomotor theories more broadly relate sickness to specific eye-movement patterns, such as nystagmus or ocular torsion, evoked by the visual stimulus \cite{ebenholtzMotionSicknessOculomotor1992}.
At an evolutionary level, Treisman's toxin-defense hypothesis explains why disturbed sensory integration should give rise to nausea and emesis, but does not specify the proximate sensory conditions that trigger them \cite{treismanMotionSicknessEvolutionary1977}.
Finally, predictive-coding formulations recast sensory mismatch as a persistent prediction error between expected and observed input, offering a unifying cognitive framework for the preceding theories \cite{nurnbergerMismatchVisualVestibularInformation2021}.

\subsection{Neural Recording and Modulation Modalities}
VIMS has historically been assessed through subjective reports, but a growing body of work uses neural recordings, and more recently non-invasive brain stimulation, to characterise the user's state during VR exposure or to actively modulate it.

Electroencephalography (EEG) is the most widely used recording modality in this literature.
Its high temporal resolution supports the study of spectral power, event-related potentials, and connectivity dynamics throughout the development of sickness, and its portability makes it compatible with HMD-based and other interactive VR paradigms \cite{cortesEEGbasedExperimentVR2023a, krokosQuantifyingVRCybersickness2022}.
These strengths are offset by noise, limited spatial specificity, particularly for deep vestibular structures, and by methodological sensitivities, such as motion artifacts, interference, reference choice, and baseline selection, which constrain the interpretation of reported effects.

Functional near-infrared spectroscopy (fNIRS) is an emerging optical technique that measures cortical hemodynamic changes via near-infrared light \cite{scholkmannReviewContinuousWave2014, PresentFutureUse}.
It offers an intermediate trade-off between temporal and spatial resolution and tolerates moderate movement, making it well suited to VR paradigms such as HMD-based roller-coaster exposure, VR locomotion, or stereoscopic viewing \cite{yeoInvestigatingCorticalActivity2024, yamamuraHemodynamicVRAdaptingUsers2021, yamamuraPleasantLocomotionReducing2020, tanimuraEffectsLowResolution3D2019}.
Because it relies on hemodynamic rather than electrical signals, its temporal resolution is substantially lower than that of EEG, and its optical penetration restricts measurement to superficial cortical regions.

Beyond recording, non-invasive neurostimulation has begun to be used to modulate sickness-relevant processes during VR or visually induced motion exposure.
Reported approaches include galvanic vestibular stimulation, transcranial direct-current stimulation, transcranial alternating-current stimulation, and related techniques \cite{benelliFrequencyDependentReductionCybersickness2023}.
By directly perturbing vestibular or cortical activity, these methods provide a causal complement to the correlational evidence offered by EEG and fNIRS.
